\documentclass[11pt]{article}
\usepackage[margin=1in]{geometry}
\usepackage{amsmath,amssymb,amsthm}
\usepackage{graphicx}
\usepackage[colorlinks=true,linkcolor=blue,citecolor=blue,urlcolor=blue]{hyperref}

\newtheorem{theorem}{Theorem}
\newtheorem{lemma}{Lemma}
\newtheorem{remark}{Remark}

\newcommand{\ellone}{\ell^1}
\newcommand{\elltwo}{\ell^2}
\newcommand{\Hn}{\ell_2^n}
\newcommand{\ellinf}{\ell^\infty}
\newcommand{\cnull}{c_0}
\newcommand{\N}{\mathbb N}

\title{A weak*-to-weak continuous operator on \texorpdfstring{$\ell^1$}{ell1} with uncomplemented kernel}
\author{Agent packet for \texttt{fa\_banach\_001}}
\date{2026-06-28}

\begin{document}
\maketitle

\section*{Source and target question}

The source paper is Flemming--Hofmann, \emph{New aspects of ill-posedness
classification in Banach spaces}, arXiv:2511.06690 \cite{FH}.  After proving
that Mazur-type operators and their nonzero compositions fail weak*-to-weak
continuity, the authors write that they were unable to find weak*-to-weak
continuous operators on $\ellone$ with uncomplemented null-space.  They then ask
whether, in $\ellone$, weak*-to-weak continuity implies complementedness of the
null-space.

\begin{center}
\includegraphics[width=\textwidth]{figures/open_problem_crop.png}
\end{center}

\section*{Result}

We give a counterexample.  There is a bounded operator
\[
        A:\ellone \longrightarrow \cnull
\]
which is weak*-to-weak continuous, but $\ker A$ is not complemented in
$\ellone$.  In particular, weak* closedness of the null-space also does not
imply complementedness in $\ellone$.

\section*{Proof intuition}

Weak*-to-weak continuity of an operator $A:\ellone=\cnull^*\to Y$ is easy to
force if every scalar functional on the range pulls back to an element of
$\cnull$.  Thus we build $A$ from scalar coordinates indexed by a dense sequence
in a subspace $E\subset \cnull$; then $A^*(Y^*)\subset E\subset \cnull$ and
$\ker A=E^\perp$.

The remaining task is to choose $E$ so that $E^\perp$ is not complemented in
$\ellone$.  If $E^\perp$ were complemented, then the quotient
$\ellone/E^\perp\cong E^*$ would be a complemented subspace of $\ellone$, hence
isomorphic to $\ellone$.  We choose
\[
        E \cong \Big(\bigoplus_{n=1}^\infty \Hn\Big)_{\cnull}.
\]
Its dual is the $\ellone$-sum of the Hilbert blocks $\Hn$, which is not
isomorphic to $\ellone$ because its dual contains uniformly 1-complemented
Hilbert blocks of unbounded dimension.

\section*{Auxiliary facts}

\begin{lemma}\label{lem:block}
The space
\[
        E_0=\Big(\bigoplus_{n=1}^\infty \Hn\Big)_{\cnull}
\]
embeds isomorphically into $\cnull$, while
\[
        E_0^*=\Big(\bigoplus_{n=1}^\infty \Hn\Big)_{\ellone}
\]
is not isomorphic to $\ellone$.
\end{lemma}

\begin{proof}
For each $n$, choose a finite $1/2$-net $F_n$ in the unit sphere of
$\Hn$.  The map
\[
        J_n h=(\langle h,u\rangle)_{u\in F_n}\in \ellinf(F_n)
\]
satisfies
\[
        \frac12\|h\|_2\leq \|J_nh\|_\infty\leq \|h\|_2.
\]
The block diagonal map $\oplus_n J_n$ embeds $E_0$ into the finite-block
$\cnull$-sum $\left(\oplus_n\ellinf(F_n)\right)_{\cnull}$, which is naturally
isometric to $\cnull$.

The dual formula follows from the standard duality of $\cnull$-sums of
finite-dimensional spaces.  To see that this dual is not isomorphic to
$\ellone$, suppose otherwise.  Then its dual
\[
        \Big(\bigoplus_{n=1}^\infty \Hn\Big)_{\ellinf}
\]
would be isomorphic to $\ellinf$.  Since $\ellinf$ is injective, this would make
$\left(\oplus_n \Hn\right)_{\ellinf}$ $\lambda$-injective for some
finite $\lambda$.  Each coordinate Hilbert block $\Hn$ is
1-complemented in that space, so the injectivity constants of the spaces
$\Hn$ would be bounded independently of $n$.  This contradicts the
classical fact that the projection/injectivity constants of finite-dimensional
Hilbert spaces tend to infinity with $n$; in fact they are of order
$\sqrt n$ \cite{Grunbaum,LT}.
\end{proof}

\begin{lemma}\label{lem:comp}
Let $E\subset \cnull$ be a closed subspace such that $E^*$ is not isomorphic to
$\ellone$.  Then its annihilator
\[
        E^\perp=\{x\in \ellone:\langle x,e\rangle=0\hbox{ for all }e\in E\}
\]
is not complemented in $\ellone$.
\end{lemma}

\begin{proof}
The quotient map $\ellone=\cnull^*\to E^*$ has kernel $E^\perp$ and is onto by
Hahn--Banach, so $\ellone/E^\perp\cong E^*$.  If $E^\perp$ were complemented in
$\ellone$, then $E^*$ would be isomorphic to a complemented subspace of
$\ellone$.  By the classical Pe\l czynski theorem, every infinite-dimensional
complemented subspace of $\ellone$ is isomorphic to $\ellone$ \cite{LT}.  This
contradicts the hypothesis on $E^*$.
\end{proof}

\section*{Counterexample}

\begin{theorem}\label{thm:counterexample}
There is a bounded weak*-to-weak continuous operator $A:\ellone\to \cnull$ whose
kernel is uncomplemented in $\ellone$.
\end{theorem}

\begin{proof}
Let $E\subset \cnull$ be an isomorphic copy of
$E_0=(\oplus_n \Hn)_{\cnull}$ as in Lemma~\ref{lem:block}.  Choose a
sequence $(v_k)_{k=1}^\infty$ in the unit ball of $E$ whose linear span is dense
in $E$.  Define
\[
        A x=\big(2^{-k}\langle x,v_k\rangle\big)_{k=1}^\infty,
        \qquad x\in \ellone=\cnull^* .
\]
This is a bounded operator from $\ellone$ into $\cnull$, since
$|2^{-k}\langle x,v_k\rangle|\leq 2^{-k}\|x\|_1$.

We verify weak*-to-weak continuity.  Let $a=(a_k)\in \ellone=\cnull^*$.  Then
for $x\in\ellone$,
\[
        \langle a,Ax\rangle
        =\sum_{k=1}^\infty a_k2^{-k}\langle x,v_k\rangle
        =\left\langle x,\sum_{k=1}^\infty a_k2^{-k}v_k\right\rangle .
\]
The series $\sum_k a_k2^{-k}v_k$ converges in norm in $E\subset\cnull$.  Hence
$A^*a\in\cnull$ for every $a\in\ellone$.  Equivalently, every scalar functional
on $\cnull$ pulls back to a weak*-continuous functional on $\ellone$.  Therefore
$A$ sends weak*-convergent nets, and in particular weak*-convergent sequences,
in $\ellone$ to weakly convergent nets in $\cnull$.

Finally,
\[
        \ker A
        =\{x\in\ellone:\langle x,v_k\rangle=0\hbox{ for all }k\}
        =E^\perp,
\]
because the span of $(v_k)$ is dense in $E$.  By Lemmas~\ref{lem:block} and
\ref{lem:comp}, $E^\perp$ is not complemented in $\ellone$.
\end{proof}

\section*{Verification and limitations}

The construction uses no numerical search.  The formula for $A$ directly gives
boundedness, weak*-to-weak continuity through the adjoint inclusion
$A^*(\ellone)\subset \cnull$, and $\ker A=E^\perp$.  The only external Banach
space facts used are standard: finite-dimensional Hilbert spaces have
unbounded injectivity/projection constants, $\ellinf$ is injective, and every
infinite-dimensional complemented subspace of $\ellone$ is isomorphic to
$\ellone$.

The counterexample answers the source question negatively in the sequential
sense used in the paper, and actually gives the stronger net-continuity
version.  It also refutes the intermediate question whether weak* closedness of
the null-space implies complementedness in $\ellone$.

\section*{Novelty check}

The run's lightweight indexes had no existing packet for arXiv:2511.06690.
Searches were run for exact and close phrases including:
\begin{itemize}
\item \texttt{weak*-to-weak continuity complemented ell\^{}1 null-space};
\item \texttt{weak* closed subspace ell\^{}1 complemented};
\item \texttt{weak*-closed subspace l\_1 uncomplemented}.
\end{itemize}
No direct later answer to the Flemming--Hofmann question was found.  The result
should nevertheless be reviewed as a candidate counterexample because it relies
on classical local-theory facts rather than on a cited paper explicitly
answering this 2025 question.

\section*{Human review recommendation}

Recommended review focus: check the standard local-theory step in
Lemma~\ref{lem:block}, namely that
$(\oplus_n \Hn)_{\ellone}$ is not isomorphic to $\ellone$ via the
unbounded projection constants of $\Hn$.  The operator construction and
kernel computation are otherwise direct.

\begin{thebibliography}{9}

\bibitem{FH}
J. Flemming and B. Hofmann,
\emph{New aspects of ill-posedness classification in Banach spaces},
arXiv:2511.06690, 2025.

\bibitem{Grunbaum}
B. Grunbaum,
\emph{Projection constants},
Trans. Amer. Math. Soc. 95 (1960), 451--465.

\bibitem{LT}
J. Lindenstrauss and L. Tzafriri,
\emph{Classical Banach Spaces I. Sequence Spaces},
Ergebnisse der Mathematik und ihrer Grenzgebiete, vol. 92,
Springer-Verlag, 1977.

\end{thebibliography}

\end{document}
